% Context from chapters/wavelets.tex; for mathematical reference, not compilation.
\section{Fast Wavelet Transform}\label{sec-fast-wavelet-transform}

   
\subsection{Discretization}

We now work over $\RR/\ZZ$.
 
We assume access to a discrete signal $a_J \in \RR^N$ with $N = 2^{-J}$ at a fixed scale $2^J$, whose entries are inner products with the scaling functions, i.e.
\eql{\label{eq-hyp-wav-samples}
	\foralls n\in \{0,\ldots,N-1\}, \quad a_{J,n} = \dotp{f}{\phi_{J,n}^P} \approx 2^{J/2}f(2^Jn),
}
for the underlying continuous function $f$. The approximation uses the normalization $\int\phi=1$. The equality assumes that acquisition gives exact access to $P_{V_J}f$, much as Shannon sampling assumes bandlimiting. This model is reasonable when the scaling functions $(\phi_{J,n})_n$ approximate the acquisition device's point-spread function. Preprocessing the measurements can sometimes improve agreement with~\eqref{eq-hyp-wav-samples}.

Starting from $a_J$, the discrete wavelet transform computes the coefficients
\eq{
	\foralls j \in \{J+1,J+2, \ldots, 0\}, \quad
	\foralls n \in \range{0,2^{-j}-1}, \quad
	a_{j,n} \eqdef \dotp{f}{\phi_{j,n}^P}, \qandq
	d_{j,n} \eqdef \dotp{f}{\psi_{j,n}^P}
}
in order of increasing $j$. After an approximation vector $a_j$ has been used to compute the next scale, it can be discarded; the detail vectors $d_j$ are retained.

The forward discrete wavelet transform on a bounded domain is thus the orthogonal finite-dimensional map
\eq{
	a_J \in \RR^N \longmapsto
	\enscond{ d_{j,n} }{ J<j\leq0,\ 0\leq n<2^{-j} } \cup \{a_0 \in \RR\}.
}
By orthogonality, the inverse transform is the adjoint.


Figure~\ref{fig-wavcoef-1d} shows examples of wavelet coefficients. For each scale $2^j$, there are $2^{-j}$ coefficients.

\myfigure{
\tabun{
\image{wavelets}{.9}{wavcoef-1d-all}\\
Wavelet coefficients $\{d_{j,n}\}_{j,n}$
}
\tabdeux{
\image{wavelets}{.48}{wavcoef-1d-0}&
\image{wavelets}{.48}{wavcoef-1d-1}\\
Signal $f$ & $d_{-7}$ \\
\image{wavelets}{.48}{wavcoef-1d-2}&
\image{wavelets}{.48}{wavcoef-1d-3}\\
$d_{-6}$ & $d_{-5}$ 
}
}{ 
	Wavelet coefficients. Top: the complete decomposition. Below: enlarged views of individual scales. 
}{fig-wavcoef-1d}


   
\subsection{Forward Fast Wavelet Transform (FWT)}

The algorithm applies a sequence of elementary operators
\eql{\label{eq-wavelet-step}
	\foralls j = J+1, \ldots, 0, \quad 
	(a_{j},d_{j}) = \Ww_j(a_{j-1})
	\qwhereq
	\Ww_j : \RR^{2^{-j+1}} \rightarrow \RR^{2^{-j}} \times \RR^{2^{-j}}
}
There are $|J|=\log_2(N)$ steps. Each $\Ww_j$ changes coordinates between orthonormal bases and is therefore orthogonal.

To describe the algorithm for $\Ww_j$, we introduce the filter coefficients $h,g \in \RR^\ZZ$
\eql{\label{eq-dfn-h-g}
	h_n \eqdef \frac{1}{\sqrt{2}} \dotp{ \phi(\cdot/2) }{\phi(\cdot-n)}
	\qandq
	g_n \eqdef \frac{1}{\sqrt{2}} \dotp{ \psi(\cdot/2) }{\phi(\cdot-n)}.
}
Figure~\ref{fig-haar-inner} illustrates the computation of these weights for the Haar system.

\begin{figure}[tbp]
\centering
\BookDrawing{tikz/wavelets-haar-inner-products}
\caption{Haar filter weights as inner products.}
\label{fig-haar-inner}
\end{figure}


We denote by $\downarrow_2 : \RR^K \rightarrow \RR^{K/2}$ the subsampling operator by a factor of 2, i.e.
\eq{
	u\downarrow_2 \eqdef (u_0,u_2,u_4,\ldots,u_{K-2}).
}
Throughout the following derivation, the scaling function, wavelet, and filters are real, and the filters decay sufficiently fast. The notation $\bar h_n=h_{-n}$ and $\bar g_n=g_{-n}$ denotes reflection. Complex filters would require conjugation as well.

\begin{prop}
One has
\eql{\label{eq-convol-subsample}
	\Ww_j(a_{j-1})= ( (\bar h \star a_{j-1})\downarrow_2, (\bar g \star a_{j-1})\downarrow_2 ),
}
where $\star$ denotes periodic convolutions on $\RR^{2^{-j+1}}$.
\end{prop}

\begin{proof}
We note that $\phi(\cdot/2)$ and $\psi(\cdot/2)$ belong to $V_0$, so one has the decompositions
\eql{\label{eq-phipsi-recurse-1}
	\frac{1}{\sqrt{2}} \phi(t/2) = \sum_n h_n \phi(t-n)
	\qandq
	\frac{1}{\sqrt{2}} \psi(t/2) = \sum_n g_n \phi(t-n)
}
Making the substitution $t \mapsto \frac{t-2^j p}{2^{j-1}}$ in~\eqref{eq-phipsi-recurse-1}, one obtains
\eq{
	\frac{1}{\sqrt{2}} \phi\pa{ \frac{t-2^j 